\subsection{The local identification model}
\label{sec:local}
%==========================================================================

\emph{The same model, read over the band a managed patient's marker actually visits. This
is where the identifiability question is posed, and Remark~\ref{rem:not_a_reduction} says
exactly how it relates to the model it is read from.}
%==========================================================================

\subsubsection{Formulation}

\begin{equation}
\dot m \;=\; s(t) \;-\; \theta\, g(m) \;-\; k_e\, m,
\qquad g(m) = \frac{m}{K_m + m}.
\label{eq:A}
\end{equation}

\begin{description}[leftmargin=1.4em, itemsep=1pt]
\item[$m(t)$] accumulated toxic metabolite, measured sparsely.
\item[$s(t)$] substrate load; partially known.
\item[$\theta$] the capacity of the enzymatic channel, its maximal velocity. Since most
pathogenic variants reduce the \emph{quantity} of functional enzyme rather than the turnover
number of an individual molecule, $\theta$ is essentially ``how much working enzyme the
patient has'', and is therefore constant and genotype-determined.
\item[$K_m$] the half-saturation level of that channel: a shape parameter of the clearance
curve at the plasma level, fitted here (Remark~\ref{rem:no_km_orig}).
\item[$g(m)$] the \emph{Michaelis-Menten saturation function}, $g \in [0,1)$.
\item[$k_e m$] alternative elimination, taken proportional to $m$.
\end{description}

The block diagram of this model is Figure~\ref{fig:blockA} of the paper, where the two
channels and the two selective inputs are drawn.

\subsubsection{Assumptions}

Three assumptions, stated here and warranted in Appendix~\ref{sec:supp}.
Constancy: $\theta$ is constant over the afebrile range, which is a restriction the
kinetics require, since variant PCCs are \emph{less} thermostable than wild
type~\cite{pcc_variants_kinetics} % C22
and would therefore be at their weakest when the load is at its highest.
Population parameters: $K_m$ and $k_e$ do not vary with genotype and may be
estimated once across the cohort, motivated for $K_m$ by the absence of a significant in
vitro difference between variant and wild-type enzyme~\cite{pcc_variants_kinetics} % C21
and for $k_e$ by the alternative channel not involving PCC, weakened by carnitine being
administered at patient-specific doses. Observable proxy: the measured marker
stands in for the intracellular substrate, which is a shortcut, since the true system has at
least two states with an uncertain map between them.



\paragraph{The marker is 3-hydroxypropionate.}
\label{sec:marker_orig}
Earlier drafts left $m$ as ``the plasma marker'' on the grounds that nothing in Sections 3
and 4 depends on the choice. That is true of the algebra and false of the decision the paper
is built around, because the three candidates differ in whether they are associated with the
endpoint at all. The choice is fixed here, and the reasons are decision-relevant rather than
analytical.

The endpoint. Over a three-year natural history cohort, a $50\%$ \emph{lower}
plasma 3-HP is associated with a decompensation rate ratio of $0.7$ (95\% CI $0.6$ to $0.9$,
$p<0.01$); a $50\%$ lower C3 with $0.6$ ($p=0.017$); and 2-methylcitrate shows no
association at all, rate ratio $1.0$, $p=0.84$~\cite{natural_history}. % C19
The conditioning of those ratios is worth stating explicitly because reading it the other
way inverts the conclusion. The ratio is attached to a lower concentration, so a rate ratio
below one means \emph{fewer events when the marker is lower}: 3-HP is a risk marker in the
expected direction for a toxic overflow metabolite, not a protective one.

C3 is excluded on the perturbation, and the argument is mechanistic. C3 is
propionylcarnitine, the product of the carnitine conjugation route. Carnitine
supplementation therefore \emph{raises} C3 while lowering the organic acids, so the
selective perturbation of Section~\ref{sec:selective} enters the observation with the wrong
sign on that channel: the input the whole identification rests on would appear to worsen the
patient. No arrangement of the algebra repairs that, and it disqualifies C3 whatever its
rate ratio. \emph{As corrected.} The identification does not rest on carnitine. The known
input is the dose train (Section~\ref{sec:therapychannel}), and carnitine is closed as a
substitute for it in Section~\ref{sec:whoto}. The exclusion of C3 stands on the narrower
ground Remark~\ref{sec:marker} states, that carnitine is a co-treatment given to every
patient and raises C3 while the organic acids fall. Whether C3 would be disqualified in a
record with no carnitine in it is not settled here, and the design does not need it to be.

2-MC is excluded on the endpoint, and its strength lies in the arm this paper
excludes. 2-MC is among the strongest correlates of plasma ammonia \emph{during} an
event~\cite{amd_phenotyping}, % C37
and by Remark~\ref{rem:scope} the behaviour through a decompensation is precisely where the
analysis here does not go. A marker that tracks the current crisis and a marker that
predicts the next one need not be the same marker, and this paper is about the second.

3-HP carries C3's immunity without its defect. It is not a carnitine conjugate, so
the sign argument above does not touch it; it has the association with the endpoint that
2-MC lacks; and it is the covariate the natural history study's own longitudinal hazard is
already fitted on~\cite{natural_history}, % C28
so the reformulation in Section~7 substitutes into machinery that exists rather than
proposing new machinery.

\begin{remark}[A marker panel, and whether the identifiability budget affords one]
\label{rem:panel}
The trial measures four markers together, 3-HP, 2-MC, C3 and
$n$-propionylglycine~\cite{trial_phase12}, which invites a multi-output observation model:
one latent pool, a channel-specific gain onto each marker, at a cost of one scaling
parameter per channel. Whether that is affordable is an identifiability-budget question, and
the budget answers it.

With $J$ channels observed as scaled copies of one pool, and \emph{independent} assay noise
on each, the assay contribution to $\sigma$ averages down as $J^{-1/2}$. The error in the
recorded load does not average at all: it is common to every channel, because every channel
is fitted against the same ordinate. So by \eqref{eq:crossover} the panel pays exactly in
the regime where the assay term dominates, which is where the load is well recorded, and
pays nothing where it is not. At a load recorded to $2\%$ a four-marker panel cuts $\sigma$
by $1.65$; at $10\%$ it cuts it by $1.07$ and at $20\%$ by $1.02$.

The recommendation is therefore conditional and in that order: buy the load record
first, and the panel only after. A panel bought first is bought against a noise term it
cannot touch. Fitting the four channels also costs three gains, and two of the four are
excluded above on grounds the averaging does not repair, so the honest version of the panel
is 3-HP and $n$-propionylglycine, a glycine conjugate that carnitine does not raise, with
$J = 2$ and a factor of $1.31$ at a load recorded to $2\%$, falling to $1.05$ at $10\%$.
\end{remark}

\subsubsection{The obstruction, and what kind it is}

The obstruction below is one of \emph{conditioning}, not of structure. Section~\ref{sec:regime} makes the distinction and disambiguates
the word. What follows is a statement about the range of metabolite levels a record
actually contains. An earlier heading here called it structural, which it is not.

Write $C(m) = \theta\, g(m) + k_e\, m$ for the total clearance, and let
$M = [m_{\min}, m_{\max}]$ be the range of metabolite levels a given patient actually
visits between clinic visits. Clinical management exists to keep $M$ narrow. The
obstruction is stated as a property of that range.

\paragraph{Two conventions, fixed here because everything downstream turns on them.}
$\Delta = m_{\max} - m_{\min}$ is the full width of $M$, not a half-width, and
$\bar m$ is the point the affine fit is taken at and the sample mean of the band. Together
they carry a constraint. The metabolite cannot be negative, so $m_{\min} \ge 0$ and
\begin{equation}
\frac{\Delta}{\bar m} \;\le\; 2.
\label{eq:bandceiling}
\end{equation}
A band wider than twice the operating point does not exist. Design figures quoted
as band widths above $200\%$ of $\bar m$ describe an interval whose foot is below zero, and
Section~\ref{sec:design} is where that matters.

The convention also moves a number. The gap between the tangent at $\bar m$ and a
least-squares fit over the band is $0.09\%$ at $\Delta/\bar m = 10\%$ and $8.8\%$ at
$\Delta/\bar m = 100\%$, where the latter means $m \in [\bar m/2,\, 3\bar m/2]$. Read as a
half-width the same phrase would mean $m \in [0,\, 2\bar m]$ and the gap would be $35.4\%$,
four times larger. The distinction is negligible in the narrow band the next proposition is
about and is not negligible in the wide band it recommends buying.

\begin{proposition}[Degeneracy]
\label{prop:degeneracy}
Suppose that over the observed range $M$ the clearance is indistinguishable from an affine
function at the resolution of the assay: there exist constants $a$ and $b$ with
\begin{equation}
\sup_{m \in M} \bigl| C(m) - (a + b\,m) \bigr| \;\le\; \varepsilon,
\label{eq:affine}
\end{equation}
$\varepsilon$ being the measurement error. ($a$ and $b$ below are the tangent coefficients
at $\bar m$; the preceding paragraph gives what a least-squares fit returns instead.)
Then data collected while the patient remains in $M$ determine $a$ and $b$, and nothing
beyond them. With $\bar m \in M$ the point at which
the fit is taken, the two coefficients are
\begin{equation}
a \;=\; \theta\,\bigl(g(\bar m) - \bar m\, g'(\bar m)\bigr) \;=\; \theta\, g(\bar m)^2,
\label{eq:intercept}
\end{equation}
\begin{equation}
b \;=\; \theta\, g'(\bar m) \;+\; k_e.
\label{eq:slope}
\end{equation}
The slope is a single scalar carrying $\theta$ and $k_e$ with the \emph{same} sign, and no
quantity of such data separates them through it. The intercept is free of $k_e$, and given
$K_m$ it determines $\theta$.
\end{proposition}

\begin{remark}[The fibre the intercept leaves open]
\label{rem:fibre}
The intercept carries $K_m$ as well, through $g(\bar m)$, and Remark~\ref{rem:no_km_orig}
forbids importing $K_m$ from in vitro work. So $(a,b)$ known \emph{exactly} still leaves a
one-parameter family. Inverting \eqref{eq:intercept} and \eqref{eq:slope},
\begin{equation}
\theta(K_m) = \frac{a\,(K_m+\bar m)^2}{\bar m^{2}},
\qquad
k_e(K_m) = b - \frac{a\,K_m}{\bar m^{2}},
\qquad 0 < K_m \le \frac{b\,\bar m^{2}}{a},
\label{eq:fibre0}
\end{equation}
the upper bound being where $k_e$ reaches zero.

How wide the fibre is, in closed form and in observable quantities. Evaluating
\eqref{eq:fibre0} at its two endpoints, the span of $\theta$ is
\begin{equation}
\frac{\theta(K_{\max})}{\theta(0^+)}
\;=\; \left(1 + \frac{b\,\bar m}{a}\right)^{\!2}
\;=\; \left(\frac{C(\bar m)}{a}\right)^{\!2},
\label{eq:span}
\end{equation}
the second equality because the tangent satisfies $a + b\,\bar m = C(\bar m)$. At steady
state $C(\bar m)$ is the \emph{recorded load} and $a$ is the \emph{fitted intercept}, so
both are observable and the span is computable for a given patient without knowing a
single parameter. It is the square of the ratio of the load the patient carries to the part
of the clearance curve that extrapolates back to the origin, so it is large exactly when the
operating point sits far down the saturation curve, which is where management puts it.

At the illustrative values this is $(322.7/24.79)^2 = 169$, which is the factor this paper
used to quote as an illustration. It is not an illustration: it is a per-patient diagnostic,
and a clinician with a fitted band and a diet record can compute it.
Equation~\eqref{eq:thetaerr} below is the noise on top of that, and is itself conditional on
$g(\bar m)$, so it is a statement about $\theta$ \emph{given} $K_m$. Nothing in Section~3
supplies $K_m$. The population assumption above grants only that it is shared across the
cohort and
estimable once, which says where it could come from without putting it in
hand, and Remark~\ref{rem:no_km_orig} forbids importing it from in vitro work. The next paragraph
is where it comes from.

One extra datum closes the fibre outright. If $k_e$ is known, from a null genotype as in
Proposition~\ref{prop:anchor}, then $K_m = (b - k_e)\bar m^{2}/a$ and $\theta$ follows from
\eqref{eq:fibre0}, exactly and with no second steady state required. The band that resolves
$(a,b)$ does the work that Proposition~\ref{prop:anchor}'s second anchor was introduced for.
\end{remark}

\begin{proof}[Argument]
The data here are steady states, where $s = C(m)$ pointwise, so under \eqref{eq:affine}
the map from $C$ to the data factors through $(a,b)$ up to a perturbation of size
$\varepsilon$, and any two parameter pairs giving the same $(a,b)$ generate observations
differing by less than the assay can resolve. For trajectories the
same conclusion needs a Gr\"onwall bound, which is not given here. The affine function in
question is the tangent to $C$ at $\bar m$, whose slope is $C'(\bar m)$ and whose intercept
is $C(\bar m) - \bar m\, C'(\bar m)$. The alternative channel contributes $k_e \bar m -
\bar m\, k_e = 0$ to the second of these, since a line through the origin has no intercept,
which leaves \eqref{eq:intercept}; the identity $g(\bar m) - \bar m g'(\bar m) = g(\bar m)^2$
is a direct computation from $g(m) = m/(K_m + m)$. Since \eqref{eq:slope} is one equation in
two unknowns, the pairs consistent with a measured $b$ alone form a line, and every point of
it is compatible with the data.
\end{proof}

Information about the residual capacity is therefore not absent from stable-patient data.
It is concentrated in the intercept, and the intercept is the one quantity a narrow band
cannot deliver. What follows conditions on $K_m$, for the reason set out in
Remark~\ref{rem:fibre}. It is the value the fitted line takes when extrapolated back to $m = 0$, and
$M$ does not contain the origin. For a least-squares affine fit to $n$ points $m_i$ with
independent errors of standard deviation $\sigma$,
\begin{equation}
\operatorname{Var}(\hat a) \;=\; \sigma^2\!\left[\frac{1}{n} +
\frac{\bar m^2}{\sum_i (m_i - \bar m)^2}\right],
\qquad
\operatorname{Var}(\hat b) \;=\; \frac{\sigma^2}{\sum_i (m_i - \bar m)^2}.
\label{eq:varfit}
\end{equation}
Both scale as the reciprocal of $\sum_i (m_i - \bar m)^2$, and their ratio tends to
$\bar m^2$ as the band narrows, so the two variances carry different units and are not
comparable to each other. What is comparable across designs is the relative error in
$\hat\theta = \hat a / g(\bar m)^2$. Writing $\Delta$ for the width of the band and taking
the $n$ samples spread across it, $\sum_i (m_i - \bar m)^2 \approx n\Delta^2/12$, and for
$\Delta \ll \bar m$,
\begin{equation}
\frac{\operatorname{sd}(\hat\theta)}{\theta} \;\approx\;
\frac{\sigma}{\theta}\;\sqrt{\frac{12}{n}}\;\frac{\bar m}{\Delta}\;g(\bar m)^{-2}.
\label{eq:thetaerr}
\end{equation}
Three factors multiply the assay's own relative error. The sample count enters as
$n^{-1/2}$, the extrapolation from the band back to the origin as $\bar m/\Delta$, and the
patient's position on the saturation curve as $g(\bar m)^{-2}$. The last two are the two
factors of the determinant $\det M = g(\bar m)^2 (m' - m)$ of
Proposition~\ref{prop:illcond}, entering inversely, so what makes that determinant small
makes $\hat\theta$ imprecise, by the same mechanism.

The two design variables in \eqref{eq:thetaerr} both appear through
$\sum_i (m_i - \bar m)^2 \approx n\Delta^2/12$, and they are not substitutes. At a fixed
band width, doubling the number of visits buys a factor of $\sqrt2$ and leaves $\Delta$
where it was; at a fixed number of visits, doubling the width of the band buys a factor of
$2$, and doubling it again buys another. Adding visits inside a fixed band does not widen
the band, which is why the two cannot be traded against each other. That distinction is the
clinical content of the result. What bounds $n$ is the cost of a clinic visit, what bounds
$\Delta$ is the safety constraint on how far a patient may be allowed to excurse, and the
factor that carries the estimate of $\theta$ is the one the safety constraint bounds.

The clinical consequence is direct. Two patients with substantially different enzymatic
capacity, compensating with different degrees of alternative clearance, share a slope, and
what tells them apart is an intercept the band resolves poorly. They are nonetheless
different patients, and any intervention that acts on one channel and not the other will
separate them.

What breaks the degeneracy is curvature. The alternative
channel $k_e m$ is a straight line: its slope is $k_e$ at every level of $m$. The enzymatic
channel bends, because $g$ saturates, so its slope $\theta g'(m)$ falls as $m$ rises. A
straight line and a bending curve can be told apart only by watching the slope change. If
the slope of $C$ is the same at the top of $M$ as at the bottom, to within what the assay
resolves, then the data are consistent with any straight line of that slope, and
\eqref{eq:intercept} and \eqref{eq:slope} are all they contain. The same curvature is what
puts $\theta$ into the intercept. A channel that runs straight through the origin
contributes nothing there, so $a$ measures the departure of the enzymatic channel from such
a line, and it vanishes with the bending that produces it.

So the two channels separate exactly when $\theta g'(m)$ changes measurably across $M$,
which happens when the enzyme is appreciably nearer saturation at $m_{\max}$ than at
$m_{\min}$. A narrow $M$ cannot deliver that: over a short interval every smooth curve looks
straight. Widening $M$ is the only route, and it means letting the metabolite climb, which
is the beginning of a decompensation. The information the identification needs and the excursion clinical
management exists to prevent are the same thing. Condition \eqref{eq:affine} is a
property of a particular patient's record, and it can be checked against that record
directly.

\begin{remark}[The scale on which this is stated]
\label{rem:no_km_orig}
Proposition~\ref{prop:degeneracy} is conditioned on the observed clearance over the
observed range. The alternative would be to compare $m$ with $K_m$ and conclude that
patients sit far below saturation, and that comparison is unavailable here.

$K_m$ in \eqref{eq:mm} is the Michaelis constant of PCC for its own substrate,
intracellular mitochondrial propionyl-CoA, reported as $0.29$ mM in purified human liver
enzyme~\cite{pcc_liver_kinetics}. The state $m$ is a plasma marker at micromolar
concentrations~\cite{amd_phenotyping}. % C23
These are different molecules in different compartments, so the inequality has no common
scale to be evaluated on and the gap between the two numbers measures nothing.

Hypothesis \eqref{eq:affine} requires no view about where $K_m$ lies. What it forgoes is
the mechanistic reading. Writing the enzymatic channel as $\theta g(m)$, with a plasma
marker in place of the enzyme's substrate, presupposes a transfer map $\psi$ from
intracellular propionyl-CoA to the marker under which $g$ keeps its Michaelis-Menten form.
This paper does not establish one, so $K_m$ here is a fitted shape parameter of the
plasma-level clearance curve.

$\thc^{\mathrm{A}}$ in \eqref{eq:thcrit} is computed through $g(\mtol)$ and inherits the same
dependence. Establishing $\psi$, or bounding how far the conclusions move across a family
of candidate maps, is a prerequisite for reading either mechanistically.
\end{remark}

\subsubsection{Perturbation as a selective input}
\label{sec:selective}

The degeneracy of \eqref{eq:slope} is broken by moving the patient, not by measuring them
more often. Two perturbations are selective, in the sense that each acts on one channel
alone.

\begin{description}[leftmargin=1.4em, itemsep=2pt]
\item[Carnitine.] Supplementation increases the capacity of the conjugation route and
therefore acts on $k_e$, leaving $\theta$ untouched. Writing the increment as
$\kappa\, u_c(t)$ with $u_c$ the recorded dose,
\begin{equation}
k_e(t) \;=\; \underbrace{k_e}_{\text{unknown}} \;+\;
\underbrace{\kappa\, u_c(t)}_{\text{known up to }\kappa}.
\label{eq:input_carn}
\end{equation}
Carnitine is administered to \emph{every} patient in the accessible cohort and its dose is
recorded in the chart.
\item[mRNA therapy.] A dose supplies enzyme and therefore acts on $\theta$, leaving $k_e$
untouched:
\begin{equation}
\theta(t) \;=\; \underbrace{\theta_{\mathrm{base}}}_{\text{unknown}}
\;+\; \underbrace{D \sum_i e^{-\lambda (t - t_i)} \mathbf{1}_{\{t \ge t_i\}}}_{\text{known: } \Delta(t)}.
\label{eq:input}
\end{equation}
This input is investigational, and when this was written no accessible patient received
it. \emph{Superseded.} It is administered, on the published phase 1/2 biweekly schedule,
and it is the input the design of Section~\ref{sec:designsec} uses. What survives of the
sentence is a statement about reach rather than about the experiment: the design is
available to patients on the therapy and to no others (Section~\ref{sec:whoto}).
\end{description}

Write the equilibrium condition before and after the perturbation, with steady levels $m$
and $m'$ respectively. Under carnitine,

\begin{equation}
\begin{aligned}
s &= \theta\, g(m) \;+\; k_e\, m, \\
s &= \theta\, g(m') \;+\; k_e\, m' \;+\; \kappa u_c\, m',
\end{aligned}
\label{eq:twoeq_carn}
\end{equation}

and under therapy,

\begin{equation}
\begin{aligned}
s &= \theta\, g(m) \;+\; k_e\, m, \\
s &= \theta\, g(m') \;+\; k_e\, m' \;+\; \Delta\, g(m').
\end{aligned}
\label{eq:twoeq}
\end{equation}

In both cases these are two linear equations in the unknown pair $(\theta, k_e)$, and in
both cases the coefficient matrix is the same,
\begin{equation}
M \;=\;
\begin{pmatrix} g(m) & m \\[2pt] g(m') & m' \end{pmatrix},
\qquad
M \begin{pmatrix} \theta \\ k_e \end{pmatrix}
= \begin{pmatrix} s \\ s - \delta \end{pmatrix},
\label{eq:system}
\end{equation}
with $\delta = \kappa u_c\, m'$ in the first case and $\delta = \Delta\, g(m')$ in the
second. The perturbation enters the right-hand side only. Hence

\begin{equation}
\det M \;=\; g(m)\,m' - g(m')\,m
\;=\; \frac{m\,m'\,(m' - m)}{(K_m + m)(K_m + m')}.
\label{eq:det}
\end{equation}

\begin{proposition}[The determinant depends on the excursion alone]
\label{prop:channel}
Let the system be moved from one steady state to another by any known perturbation: an
increment on either channel, a change in the load, or any combination of these. Let $m$ and
$m'$ be the steady metabolite levels before and after. Then in every case the
$2 \times 2$ system for $(\theta, k_e)$ has the coefficient matrix $M$ of
\eqref{eq:system} and the determinant \eqref{eq:det}, which depends on $m$, $m'$ and $K_m$
alone. The system is two equations in two unknowns given $K_m$; with $K_m$ unknown it is two
equations in three, and Remark~\ref{rem:fibre} is where that is resolved. The perturbation enters only the right-hand side, as
$\delta = \kappa u_c\, m'$ for carnitine, $\delta = \Delta\, g(m')$ for the therapy, and
$\delta = s - s'$ for a change of load.
\end{proposition}

\begin{remark}[How far channel-blindness carries, and where it stops]
\label{rem:channel_scope}
Proposition~\ref{prop:channel} is proven at steady state and for the pair $(\theta, k_e)$.
Both restrictions have been checked against the trajectory, because the use the paper makes
of the proposition -- the question posed on carnitine and answered for the therapy --
quietly drops them.

The steady-state restriction is harmless. Integrating the model with its relative
sensitivities and comparing the observability Gramians of two inputs chosen to produce the
\emph{same} trajectory, the $(\theta, k_e)$ sub-Gramians differ by at most $5\%$ in
determinant and $2.3\%$ in $\operatorname{sd}(\ln\hat\theta)$, over excursions from $5\%$ to
$76\%$. The mechanism is exact rather than numerical. The two sensitivity forcings,
$-\theta g(m)$ and $-k_e m$, are functions of the state alone and are therefore identically
channel-blind; the only channel dependence anywhere is the decay rate they share, and at a
matched trajectory
\begin{equation}
J_{\Delta}(t) - J_{u_c}(t) \;=\; -\,\kappa u_c(t)\, g(m(t)) \;\le\; 0.
\label{eq:jgap}
\end{equation}
The therapeutic input decays the sensitivities more slowly, so it returns the larger
Gramian, always. Posing the question on carnitine is therefore \emph{conservative}: it
understates what the therapy would deliver, by a few per cent.

The two-parameter restriction is not harmless, and it is the one the paper relies
on. Adding $K_m$ to the parameter vector -- which is the situation Remark~\ref{rem:fibre}
describes, since the fibre exists precisely because $K_m$ is unknown -- the two \emph{input
routes} stop being informationally substitutable by six orders of magnitude in $\det G$ and
by a factor of $46$ to $807$ in $\operatorname{sd}(\ln\hat\theta)$. (That is carnitine
against the therapy as ways of buying information, and it is a statement about
conditioning. It is not the channel exchange of Proposition~\ref{prop:exchange}, which is
an exact parameter symmetry and which this model does not have.) The reason is in the form of the forcing, not
in any constant. The $K_m$ sensitivity
is forced by $-(\theta + \Delta(t))\,K_m\,\partial g/\partial K_m$, and that prefactor is
modulated by the therapeutic input and constant under carnitine. The ratio of the $K_m$
forcing to the $\theta$ forcing is $-(1 + \Delta(t)/\theta)\,K_m/(K_m + m)$: the therapy
moves it at will, and carnitine can move it only through $m$. At an operating point with
$\bar m/K_m = 0.1$ the two sensitivities are collinear to $|\!\cos| = 0.99994$ under
carnitine and $0.967$ under the therapy.

So the excursion is channel-blind and the curvature is not. Carnitine buys the
excursion, which is what Proposition~\ref{prop:channel} says and what is true. It cannot buy
the separation of $\theta$ from $K_m$ except through the state, and it buys that only by
carrying the patient into the neighbourhood of $K_m$: moving the operating point from
$m/K_m = 0.1$ to $m/K_m = 2$ drops the collinearity to $0.982$ and raises $\det G$ by
fifteen orders of magnitude. That excursion is the one clinical management exists to
prevent, which is the obstruction of Section~\ref{sec:structure} arriving by a second route.
\end{remark}

\begin{proposition}[Formal identifiability, practical ill-conditioning]
\label{prop:illcond}
\emph{Given $K_m$ and the input gain}, the determinant \eqref{eq:det} is non-zero whenever
$m \neq m'$, so the pair $(\theta, k_e)$ is identifiable from steady-state data before and
after any such perturbation. The second proviso is not free: the right-hand side of \eqref{eq:system} is $\delta = \kappa u_c m'$ for carnitine,
and $\kappa$ is never assigned a value anywhere in this paper (Table~\ref{tab:paramsfull}), so on
that channel the system is two equations in \emph{three} unknowns before $K_m$ is even
raised. On the therapeutic channel $\delta = \Delta g(m')$ and $\Delta$ is known from the
dosing schedule, which is a second respect, alongside
Remark~\ref{rem:channel_scope}, in which the two channels are not interchangeable. It factorises exactly, with no fit and no hypothesis:
\begin{equation}
\det M \;=\; g(m)\,m' - g(m')\,m \;=\; g(m)\,g(m')\,(m' - m).
\label{eq:detexact}
\end{equation}
Over a band narrow enough for \eqref{eq:affine} the two saturations coalesce, and
\begin{equation}
\det M \;\approx\; g(\bar m)^2\,(m' - m),
\label{eq:detgamma}
\end{equation}
the geometric mean read at the midpoint. Comparing with \eqref{eq:intercept},
$g(\bar m)^2 = a/\theta$: the two propositions are about the same object, one fitted to the
total clearance and one to the saturation function alone.
\end{proposition}

The approximation in \eqref{eq:detgamma} is not free. At $m = 8$ and $m' = 14$ with
$K_m = 100$, a $55\%$ excursion, it overstates \eqref{eq:detexact} by $8\%$, so the
narrow-band form flatters the conditioning exactly where the paper recommends widening the
band.

The determinant is small for two independent reasons, which are different clinical
situations. It is proportional to the excursion $m' - m$, so a patient held in a narrow band
gives a small determinant however far up the saturation curve they sit, and proportional to
$g(m)g(m')$, so a patient low on the curve gives a small determinant however wide the
excursion. Only a patient who is
both excursing and appreciably saturated escapes both, and that patient is in the middle of a
decompensation.

Proposition~\ref{prop:channel} is what made the analysis possible when the therapy was
prospective. The identifiability question could be posed and answered on the perturbation
actually present in the records, namely carnitine, and the answer transferred to the
perturbation that was not, namely the therapy. What has to be earned in each case is the
excursion $m \to m'$, and the channel through which it is bought is immaterial to the
information obtained. \emph{Two corrections.} The transfer holds at fixed $K_m$ and not
otherwise: with $K_m$ in the parameter vector the two routes separate by six orders of
magnitude in $\det G$ (Remark~\ref{rem:channel_scope}). And the detour is no longer
needed, because the therapy is administered and the question is posed on it directly
(Section~\ref{sec:therapychannel}).

Figure~\ref{fig:A3} carries the arithmetic. Two patients with threefold different
enzymatic capacity, indistinguishable before the perturbation, are driven apart by a
carnitine increment and by an enzymatic increment in turn, and the separation achieved is
plotted against the excursion in $m$ that each buys. The two curves coincide, which is
Proposition~\ref{prop:channel} in numbers, and both stay near $1\%$, which is well under
the analytical and biological variability of the assay.

\subsubsection{What the accessible input costs}
\label{sec:instruments}

\begin{remark}[Why the carnitine route is priced here, and why the paper does not take it]
\label{rem:instruments_local}
This subsection prices the design on carnitine, which is the route a reader will try, and
it is not the route the paper takes. The therapy is administered, so the instrument
question is not reached, and it would not help if it were: carnitine acts on $k_e$, so it
cannot move the first factor of $-(1 + \Delta(t)/\theta)\,\varphi(m)$ at all, and on the
footing that gives the dose train $\operatorname{sd}(\ln\hat\theta) = \num{sd_monthly}$, adding
recorded carnitine steps gives $1.91$ to $2.02$, \emph{worse} than leaving them out
(Section~\ref{sec:whoto}). The exogeneity an instrument buys is not the binding
constraint. What survives below is the structural half of the argument, and the directions:
the instruments move the patient the wrong way along the axis that carries the information,
which is the same finding Section~\ref{sec:whoto} reaches by a different route.
\end{remark}

Proposition~\ref{prop:channel} equalises the two channels informationally, given $K_m$. It
does not equalise them \emph{statistically}, and that was the central compromise of this
work: the carnitine dose is adjusted in response to the patient's state, so it is confounded
by indication, where a therapeutic schedule fixed in advance is exogenous by construction
and was, when this was written, administered to nobody accessible. The last clause no longer
holds, and with it the compromise: the therapeutic schedule is the one the design uses.

The way through is to use only the quasi-exogenous part of the variation in $u_c$. Three
such instruments are present in routine care, and their directions are what decides them:
weight-based dose recalculation after growth raises the dose
and \emph{lowers} $m$; cyclical intestinal antibiotics lower the load and lower $m$; only a
supply interruption raises $m$, and its amplitude is set by a stock-out rather than by a
design. A change of load carries the same determinant as a channel perturbation and costs
more, needing both loads known where the channel route needs one. None is randomised.
Appendix~\ref{sec:supp} gives each in full, and Table~\ref{tab:excursion} prices them.

\subsubsection{What we would like to establish}

\begin{enumerate}[leftmargin=*, itemsep=2pt]
\item Conditions under which a perturbation schedule, treated as a design variable,
      yields acceptable conditioning without approaching the safety constraint. By
      Proposition~\ref{prop:channel} the design problem is the same, \emph{given $K_m$},
      whether the schedule is one of carnitine doses or of therapeutic doses. It is not the
      same with $K_m$ free (Remark~\ref{rem:channel_scope}), and the item is answered for
      the therapeutic schedule in Section~\ref{sec:designsec}. What such a schedule has to buy is
      already fixed by \eqref{eq:intercept} and \eqref{eq:thetaerr}. The residual capacity
      is carried by the intercept, and the intercept is bought by the width of the band,
      where adding visits inside a fixed band does not widen it, so a schedule is to be
      judged on the excursion it produces.
\item Whether the decision-relevant quantity is better conditioned than the parameters.
      Writing the critical capacity from the equilibrium condition at peak load and
      tolerable level,
      \begin{equation}
      \spk = \thc^{\mathrm{A}}\, g(\mtol) + k_e \mtol
      \qquad \Longrightarrow \qquad
      \thc^{\mathrm{A}} = \frac{\spk - k_e \mtol}{g(\mtol)},
      \label{eq:thcrit}
      \end{equation}
      the safety margin is
      \begin{equation}
      \mu \;=\; \theta - \thc^{\mathrm{A}} \;=\; \theta \;-\; \frac{\spk}{g(\mtol)}
      \;+\; \frac{k_e \mtol}{g(\mtol)}.
      \label{eq:margin}
      \end{equation}
      Both parameters enter \eqref{eq:margin} with a \emph{positive} coefficient, and
      \eqref{eq:slope} carries them positively too, so moving along the direction the data
      cannot resolve raises one and lowers the other and their contributions to $\mu$
      partially cancel. Parametrising that direction by $K_m$, the fitted pair $(a,b)$ of
      \eqref{eq:intercept} and \eqref{eq:slope} gives
      \begin{equation}
      \theta(K_m) = \frac{a\,(K_m + \bar m)^2}{\bar m^{2}},
      \qquad
      k_e(K_m) = b - \frac{a\,K_m}{\bar m^{2}},
      \label{eq:fibre}
      \end{equation}
      and substituting into \eqref{eq:margin}, with $P = \spk - b\,\mtol$,
      \begin{equation}
      \mu(K_m) \;=\; K_m\left[\frac{2a}{\bar m} - \frac{P}{\mtol}
      - \frac{a\,\mtol}{\bar m^{2}}\right] \;+\; (a - P).
      \label{eq:mufibre}
      \end{equation}
      The $K_m^2$ terms cancel identically. The margin is \emph{linear} along the
      unidentified direction where $\theta$ is quadratic in it, so the decision quantity is
      better conditioned than the parameter it is built from, by one order. Its whole
      sensitivity to the unresolved direction is the single coefficient in
      \eqref{eq:mufibre}, computable from the fitted $(a,b)$ without resolving $K_m$.
      This is a change of order and not a claim of good conditioning: $\mu$ still varies
      several-fold over the admissible range of $K_m$ (Figure~\ref{fig:A7}). What survives
      is the \emph{sign}, and because $\mu$ is linear the sign is constant along the whole
      fibre unless the two endpoints disagree. They disagree only for $\spk$ inside
      \begin{equation}
      \spk \in \left(\;\frac{K_{\max} S + a}{1 + K_{\max}/\mtol} + b\,\mtol,
      \;\; a + b\,\mtol \;\right),
      \qquad S = \frac{2a}{\bar m} - \frac{a\,\mtol}{\bar m^{2}},
      \label{eq:window}
      \end{equation}
      with $K_{\max} = b\bar m^{2}/a$, the two roots of $\mu(K_{\max}) = 0$ and
      $\mu(0) = 0$. Outside that window the decision is determined by the data even
      though $\theta$ is not, and at the illustrative values $\theta$ is unknown there to a
      factor of $169$. Inside it no width of band settles the decision, because every point
      of the fibre fits the data equally well and they do not agree on the sign.

      The window is narrow where it matters. At $\mtol = 40$ it is
      $(3.25, 3.77)$ times the patient's baseline load, $14.8\%$ wide about its centre,
      and an illness of the amplitude the equilibrium reduction draws in Figures~\ref{fig:C1} and
      \ref{fig:C2}, half again the baseline load, sits well below it with $\mu > 0$ at
      both ends. Its width is set by the tolerable level and closes as that convention
      approaches the band the patient is actually observed in: $1.2\%$ at $\mtol = 15$,
      $14.8\%$ at $40$, $67\%$ at $160$. This replaces a tolerance on $\theta$, which the
      note never states, with a question about $\spk$ and $\mtol$, which it can compute
      from $(a, b)$ per patient.

      Two patients matched at $\bar m$ and sharing $K_m$ separate exactly:
      \begin{equation}
      \mu_A - \mu_B \;=\; (k_{e,A} - k_{e,B})\,(\mtol - \bar m),
      \label{eq:musplit}
      \end{equation}
      since point-matching forces
      $\theta_A - \theta_B = -(k_{e,A} - k_{e,B})(K_m + \bar m)$. Two observationally
      identical patients diverge linearly in how far the tolerable level sits above the
      operating point, and coincide when $\mtol = \bar m$.
\item \label{sec:invariance}Whether the decision $\theta \gtrless \thc$ is invariant under
      a common strictly increasing reparametrisation of the activity scale. It is, provided
      $\thc$ is \emph{estimated empirically on the same distorted scale} rather than
      computed from \eqref{eq:thcrit}, so the requirement is assay \emph{consistency} across
      variants and not absolute calibration. It fails if the distortion differs by mechanism
      class, and it fails whenever $\thc$ is computed rather than estimated, because
      \eqref{eq:thcrit} is not equivariant. Appendix~\ref{sec:supp} gives the argument, and
      records that this pulls against Proposition~\ref{prop:intrinsic}: an intrinsic
      threshold is a structural location, and a monotone rescaling of the activity axis does
      not preserve it. Well-posedness and invariance are bought from different accounts.
\end{enumerate}

\begin{remark}[Which threshold the cancellation is proven against, and what that costs]
\label{rem:which_threshold}
The cancellation in \eqref{eq:mufibre} is derived against $\thc^{\mathrm{A}}$ of
\eqref{eq:thcrit}, which imports $\mtol$, a clinical convention. Proposition
\ref{prop:intrinsic} defines a threshold that does not. The natural question is whether the
cancellation survives the substitution, and it does not, for a reason worth stating
precisely: the premise is absent, not the conclusion false.

The cancellation is a statement about a degeneracy. It says that along the one-parameter
family a band leaves open, the margin is linear where $\theta$ is quadratic. In the model
Proposition~\ref{prop:intrinsic} lives in, that family does not exist. The alternative
channel there is curved, so it contributes to \emph{both} coefficients of an affine fit, and
subtracting its contribution leaves an enzymatic pair $(a_e, b_e)$ from which
$K_m = \bar m^2 b_e/a_e$ follows outright. Given $\alpha$, a fitted tangent determines
$(\theta, K_m)$ uniquely, to machine precision at the illustrative constants. There is no
fibre for a margin to be linear along.

So the margin result is stated here for the operational threshold, and that is a limitation
rather than a choice. Two things go with it.

First, the identity $a = \theta g(\bar m)^2$ is a property of this model and not of
the one it is a reduction of. It holds when the alternative channel is linear through the
origin. Under the parent model the intercept is
$a = \theta g(\bar m)^2 + [\alpha(\bar m) - \bar m\,\alpha'(\bar m)]$, and at the nominal
constants the second term is $374$ against the first's $61$. The alternative channel owns
the intercept there, which is exactly the term the capacity was supposed to be carried by.

Second, the intrinsic threshold does not eliminate the convention; it relocates
it. What leaves is $\mtol$, an analytical tolerance with nothing behind it. What arrives is
$\alpha$'s constants: $K_c$ and $M_i$, which nothing measures, $\rho_0$, and above all the
calibration that fixes $R_0$, which is the claim that a patient with no residual enzyme
cannot be maintained at baseline load. The relocation is a real gain, because that last one
is a clinical assertion a clinician can dispute, where $\mtol$ is a number chosen to make a
formula close. It is not an elimination, and Remark~\ref{rem:shape} prices what arrives.
\end{remark}

\subsubsection{What the design would have to deliver}
\label{sec:design}

\begin{remark}[This design is the carnitine one, not the paper's]
\label{rem:design_carnitine}
Everything in this subsection prices a design whose $n$ counts quasi-exogenous carnitine
steps. That is not the design the paper specifies. The dose train supplies the excitation,
$n$ counts outpatient draws rather than instrument events, and twenty-five monthly draws
across dose cycles over two years reach $\operatorname{sd}(\ln\hat\theta) = \num{sd_monthly}$
(Section~\ref{sec:designsec}). The negative answer below is therefore the answer for the
carnitine route, and it is the reason the therapy is the design's condition of existence
rather than an improvement on an otherwise workable record (Section~\ref{sec:whoto}). The
instrument arithmetic below, a floor of $53\%$ at six instrument events and the wrong
direction of travel, is what closes that route.
\end{remark}

Equation \eqref{eq:thetaerr} is derived above and left unevaluated. Evaluating it
answers the question of Table~\ref{tab:models} \emph{on the carnitine route}, and the
answer is negative.

\paragraph{The relative precision is identified, though the parameter is not.}
$\hat\theta = \hat a/g(\bar m)^2$ is the fitted intercept times a constant, so
$\operatorname{sd}(\hat\theta)/\theta = \operatorname{sd}(\hat a)/a$ exactly and the fibre
coordinate cancels. The $g(\bar m)^{-2}$ of \eqref{eq:thetaerr} is not a separate unknown:
since $a = \theta g(\bar m)^2$, the product $(\sigma/\theta)g(\bar m)^{-2}$ is $\sigma/a$,
and $a$ is fitted. So
\begin{equation}
\frac{\operatorname{sd}(\hat\theta)}{\theta} \;\approx\;
\frac{\sigma}{a}\,\sqrt{\frac{12}{n}}\;\frac{\bar m}{\Delta},
\label{eq:relerr}
\end{equation}
every term of which is observable. Remark~\ref{rem:fibre} leaves $\theta$ unknown to a
factor of $169$; it leaves its \emph{relative} precision known. The two are different
things and they multiply. \eqref{eq:relerr} says how tightly the data pin $\theta$ down
\emph{relative to wherever it sits}, and Remark~\ref{rem:fibre} says the paper does not know
where that is. Neither closes the other, and the fibre remains the only unquantified
uncertainty in the location.

\paragraph{$\sigma$ and $n$.} The ordinate of the fit is $C = s - \dot m$, and
neither term is measured. At the one documented schedule, six-monthly visits
(Table~\ref{tab:epistemic}), $\dot m$ is not recoverable at all: the metabolite's time
constant is days. The accessible design is therefore steady states, where $\dot m = 0$ and
$C(m_i) = s_i$. The steady-state framing is what makes the argument for
Proposition~\ref{prop:degeneracy} immediate. And $n$ counts \emph{instrument events}, the
quasi-exogenous carnitine steps of Section~\ref{sec:instruments}, not clinic visits, which
is a much smaller number. On the design the paper specifies it counts draws, because the
dose train and not an instrument supplies the perturbation, and a draw is then an
observation of a known input (Remark~\ref{rem:design_carnitine}).

\paragraph{Two noise terms, and which one binds.} At steady state the ordinate is the
recorded load and the abscissa is the assayed metabolite, so an abscissa error $\sigma_m$
acts on the fit like an ordinate error $b\,\sigma_m$ and the two add in quadrature. The two
terms balance when
\begin{equation}
b\,\mathrm{CV}\,\bar m \;=\; f\, C(\bar m),
\qquad\text{i.e.}\qquad
\mathrm{CV} \;=\; \frac{C(\bar m)}{b\,\bar m}\, f \;\approx\; 1.08\, f,
\label{eq:crossover}
\end{equation}
with $f$ the relative error in the recorded load, because $b\bar m$ and $C(\bar m)$ differ
by under $8\%$ at these values. The assay term therefore dominates at every load
recorded to better than roughly the assay's own coefficient of variation. An earlier
version of this paragraph asserted the reverse without checking it. At a load recorded to
$2\%$ and a $5\%$ assay the assay term is $14.9$ against the load's $6.5$, and the total
$\sigma$ is two and a half times what the load alone would give.

\paragraph{The variance has a floor, and \eqref{eq:thetaerr} dropped it.} Equation
\eqref{eq:varfit} has two terms and \eqref{eq:thetaerr} keeps only the second. The first is
$1/n$, and it does not move with the band width, so it binds in exactly the wide-band regime
this section recommends buying. Writing $S = \sum_i (m_i - \bar m)^2$,
\begin{equation}
\frac{\operatorname{sd}(\hat\theta)}{\theta} \;=\; \frac{\sigma}{a}
\sqrt{\frac{1}{n} + \frac{\bar m^2}{S}}
\;\;\ge\;\; \frac{c(n)\,\sigma}{a\sqrt n},
\qquad
c(n) = \sqrt{1 + \frac{3(n-1)}{n+1}} \;\le\; 2,
\label{eq:floor}
\end{equation}
for $n$ evenly spaced samples, the bound following from \eqref{eq:bandceiling} and attained
only as the band's foot reaches $m = 0$. The best any design supported on $[L, H]$ can do is
$\operatorname{sd}(\hat a) = \sigma (H+L)/(\sqrt n (H-L))$, obtained by splitting the samples
between the two ends in the optimal proportion, which is $c(n) \to 1$: an evenly spaced band
is at worst a factor of two off the best use of the same excursion, and neither escapes
$\sigma/(a\sqrt n)$.

The consequence is a change of kind rather than of value. The design requirement
splits into two pieces that do not trade against each other: a floor on $n$, set by
$\sigma$ alone, below which no excursion whatever delivers the precision; and, above that
floor, a requirement on the excursion. The published version of this section had no floor
in it, and so could quote band widths of $460\%$, $920\%$ and $1841\%$ of $\bar m$ as design
requirements. By \eqref{eq:bandceiling} none of those is a band.

\paragraph{The numbers, as a function of what the load is known to.} Nothing in this paper
bounds $f$, so nothing below is quoted at a point. The second criterion is that condition
\eqref{eq:affine} be checkable on the same records: a quadratic fit has
$\operatorname{sd}(\hat c_2)$ against a true $|c_2| = \theta K_m/(K_m+\bar m)^3$, and it is
stated here as \emph{the top of the excursion} rather than as a band width, because the
curvature of $C$ is at high $m$, so that band must run upward, and because the top of the
excursion is what a safety constraint bounds and what a clinician can be asked for. Taking
$n = 6$ instrument events and a $5\%$ assay:

\begin{center}
\scriptsize
\begin{tabular}{@{}llllllll@{}}
\toprule
\textbf{load error} & \textbf{$\sigma_{\mathrm{load}}$} & \textbf{$b\sigma_m$} & \textbf{dominant} & \textbf{floor at $n=6$} & \textbf{$n$ for $20\%$} & \textbf{band foot} & \textbf{excursion top} \\
\midrule
2\% & 6.5 & 14.9 & assay & 53\% & 43 & 0.22$\,\bar m$ at $n=86$ & 2.9$\,\bar m$ \\
5\% & 16.1 & 14.9 & load & 72\% & 79 & 0.21$\,\bar m$ at $n=158$ & 3.3$\,\bar m$ \\
10\% & 32.3 & 14.9 & load & 117\% & 206 & 0.21$\,\bar m$ at $n=412$ & 4.1$\,\bar m$ \\
20\% & 64.5 & 14.9 & load & 218\% & 714 & 0.21$\,\bar m$ at $n=1428$ & 5.4$\,\bar m$ \\
\bottomrule
\end{tabular}
% floor is 2 (sigma/a)/sqrt(n), which no band width beats; band foot is where a uniform band from m_bar must reach for 20% at twice the n in the previous column; excursion top is where it must reach for the curvature of C at two standard errors, n = 6
\end{center}

Three readings, in order of how much they change the conclusion.

\begin{enumerate}[leftmargin=*, itemsep=2pt]
\item \textbf{At six instrument events the question of $20\%$ precision does not arise.}
      The floor is $53\%$ even at a load recorded to $2\%$. Reaching $20\%$ needs $43$
      instrument events at that load and $714$ at a load recorded to $20\%$, and $n$ counts
      quasi-exogenous carnitine steps, of which a chart holds a handful.
\item \textbf{The curvature criterion asks the patient to go up, and the instruments only
      take them down.} Carnitine supplementation raises clearance and lowers $m$; so does a
      weight-based dose increase after growth, and so do the cyclical antibiotics. The one
      instrument in Section~\ref{sec:instruments} that raises $m$ is an involuntary dose
      reduction, whose amplitude is set by a stock-out rather than by a design. The
      requirement is an excursion to between $2.9$ and $5.4$ times the operating point, and
      it is not that the instruments fall short of it by a factor; it is that they move the
      patient the other way.
\item \textbf{A shortfall stated as a factor does not survive stating both sides the same
      way.} A requirement written as $\Delta/\bar m$ against an achievement written as
      $\Delta/m_{\max}$, with a $\sigma$ carrying only the load term, is not a like-for-like
      ratio, and the number it yields is an artefact of the mismatch.
\end{enumerate}

\begin{remark}[The noise is in the abscissa, and the fit announces its own failure]
\label{rem:eiv}
At steady state the recorded load is the ordinate and the assayed metabolite the abscissa,
so \eqref{eq:varfit} is the wrong variance: the errors are in $x$. Classical attenuation
pulls $\hat b$ towards zero and pushes $\hat a$ up, which alone would inflate $\hat\theta$.
What that does to the margin depends on whether the fit is constrained. Unconstrained,
$\hat k_e = \hat b - \hat a K_m/\bar m^2$ goes negative, $\thc^{\mathrm{A}}$ rises, and
$\hat\mu$ appears to change sign. That reading is an artefact: a negative $k_e$ is not a
parameter set. Imposing $k_e \ge 0$ leaves the margin biased downward by about $7\%$ at
$\sigma_m/\Delta = 20\%$ and no more, which is small beside everything else here.

What survives is better than the bias. At $\sigma_m/\Delta = 20\%$, $96\%$ of unconstrained
fits return $\hat k_e < 0$, and that is checkable from the fit alone, with no knowledge of
the truth. The design does not merely mislead, it reports that it has failed. A companion
diagnostic is available to the analyst though not to the clinic: inflating $\hat a$ shrinks
the fibre \eqref{eq:fibre0} from above, and past about $1.6$ times the true $a$ the fibre
excludes the true $\theta$ altogether.
\end{remark}

\begin{remark}[The one excursion-free route, and why it is kept]
\label{rem:anchor_kept}
The null-genotype anchor is the only construction in this paper that recovers $\theta$
without an excursion: a patient with $\theta_j = 0$ has $s_j = k_e m_j$ outright, and by
Remark~\ref{rem:fibre} a known $k_e$ collapses the fibre to a point. Everything else here
is built on excursions, so the anchor has to be either promoted or explained.

It is kept, as a remark, because it has the same shape as the paper's main
argument. Remark~\ref{rem:anchor_var} prices it, and the price inverts: the anchor
\emph{attenuates} a load error for a mild patient and \emph{amplifies} it for a severe one,
crossing one at $\theta \approx 1100$ and reaching $3.7$ at the $\theta = 300$ the
equilibrium reduction uses. So the one route that needs no excursion works for the patients
who need it least, which is the same inverse-usefulness the excitation condition has: the
information is easiest to get from the patients whose management is going worst. Retiring it
would remove an instance of the paper's own pattern; promoting it would misrepresent what it
delivers for the severe patient.
\end{remark}

\begin{remark}[The anchor's cost, a ratio that inverts for the severe patient]
\label{rem:anchor_var}
Proposition~\ref{prop:anchor} is demonstrated noise-free, which is the standard this section
spends its length arguing against. A null patient has $\theta_j = 0$, so $s_j = k_e m_j$ and
a relative error in the recorded load enters $\hat k_e$ one for one. Through
$K_m = (b - k_e)\bar m^2/a$ and \eqref{eq:fibre0} it reaches
\begin{equation}
\left|\frac{\mathrm{d}\ln K_m}{\mathrm{d}\ln k_e}\right|
= \frac{k_e (K_m + \bar m)^2}{\theta\,K_m},
\qquad
\left|\frac{\mathrm{d}\ln \theta}{\mathrm{d}\ln k_e}\right|
= \frac{2 k_e (K_m + \bar m)}{\theta}.
\label{eq:anchorfac}
\end{equation}
Both scale as $1/\theta$, so this is a ratio, the same shape as the
load-error amplification of Figure~\ref{fig:A2}. At $\theta = 3000$ they are $0.20$ and
$0.37$ and the anchor attenuates; they pass one at $\theta = k_e(K_m+\bar m)^2/K_m = 605$
and $\theta = 2k_e(K_m+\bar m) = 1100$; at the parent's own $\theta = 300$ they are $2.0$ and
$3.7$ and the anchor \emph{amplifies}. It is therefore worst for the severe patients the
note is about, and the favourable figure belongs to the mild end of the range.
\end{remark}

\subsubsection{Known defect}

Since $k_e m \to \infty$, total clearance is unbounded and every load is eventually balanced,
so the local reduction predicts that no crisis can occur. It is adequate for the
identifiability analysis, which concerns the stable patient, and inadequate for anything
concerning decompensation, which is what Section~\ref{sec:dynamic} is for.

\vspace{0.5em}
\hrule
\vspace{1em}

%==========================================================================
